\begingroup
\small
\section*{Sources and drafting record}
\addcontentsline{toc}{section}{Sources and drafting record}

The structural source for every numbered section and appendix is
Fauzan's manuscript \cite{F}, dated 17 September 2026, version 1,
32 pages. Its equations are cited by their original numbers throughout.
The title below records the title of that manuscript; it is not an
endorsement of its main claim. The downloaded PDF is preserved locally,
with SHA-256
\begin{center}\small\ttfamily
a42c05ccc46b9a7c36428913453c4617202\\
c59109713ba20f464ab9ec640d0da.
\end{center}

\begin{center}
\small
\begin{tabular}{@{}>{\raggedright\arraybackslash}p{0.18\textwidth}>{\raggedright\arraybackslash}p{0.22\textwidth}>{\raggedright\arraybackslash}p{0.50\textwidth}@{}}
\toprule
Draft & Source location & Treatment\\
\midrule
Sections 1--2 & pp.~1--5 & New overview; ninth-order functional, determinant, degree, and positive integral.\\
Section 3 & pp.~5--9 & Ninth-order pullback, completion, distribution, and small-prime calculation.\\
Section 4 & pp.~9--13 & New weights, moment cutoff, and outer costs; inner analytic transfer marked formal.\\
Section 5 & pp.~13--17 & New normalization budget, with dependencies carried forward.\\
Section 6 & pp.~17--20 & Same leading field; new weight bounds and scaling errors.\\
Section 7 & p.~21 & Sign comparison and conditional completion.\\
Appendix A & pp.~21--26 & Same measure and elementary formulas; source's 684 interval checks imported as a certificate.\\
Appendix B & pp.~26--28 & New exact rational costs and revised margins.\\
Appendix C & pp.~28--31 & New scaling and residues; conditional approximation and height calculations.\\
\bottomrule
\end{tabular}
\end{center}

The standard identities used to construct the new functional are
documented in the NIST Digital Library of Mathematical Functions:
Euler's even-zeta formula \cite{D1}, Hermite's integral \cite{D2},
the Bernoulli difference, reflection, and multiplication identities
\cite{D3}, and von Staudt--Clausen \cite{D4}.
The exact rational calculations and finite valuation examples are
calculations for this exploration; they are not external sources.
Earlier local seventh-order notes were used as working context, with
the ninth-order substitutions written out here.

Six separate section assignments used GPT-6 Astra at extra-high
reasoning (\texttt{xhigh}, not \texttt{ultra}): Section~2; Section~3;
Section~4; Section~5 with Appendix~B; Section~6 with Appendix~A; and
Section~7 with Appendix~C. The coordinating agent wrote the introduction,
aligned the shared definitions and constants, and assembled and checked
the document's layout. This division of work is not independent peer
review. Handoff notes and the shared editorial contract are included
with the editable source bundle.

\begin{thebibliography}{9}
\raggedright
\bibitem{F} Aabir Fauzan, \emph{$\zeta(5)$ is irrational},
17 September 2026, version 1, Zenodo.
\href{https://zenodo.org/records/22826419}{doi:10.5281/zenodo.22826419}.
\bibitem{D1} NIST Digital Library of Mathematical Functions,
\href{https://dlmf.nist.gov/25.6.E2}{equation 25.6.2},
even values of the Riemann zeta function.
\bibitem{D2} NIST Digital Library of Mathematical Functions,
\href{https://dlmf.nist.gov/25.11.E29}{equation 25.11.29},
Hermite integral for the Hurwitz zeta function.
\bibitem{D3} NIST Digital Library of Mathematical Functions,
\href{https://dlmf.nist.gov/24.4}{Section 24.4}, particularly
equations 24.4.1, 24.4.3, and 24.4.18.
\bibitem{D4} NIST Digital Library of Mathematical Functions,
\href{https://dlmf.nist.gov/24.10.i}{Section 24.10(i)},
von Staudt--Clausen theorem.
\end{thebibliography}
\noindent\small Source URLs last consulted for this exploration:
14 November 2026. The source bundle includes a machine-readable ledger
of locations, scope, and file hashes.
\endgroup
